% Original applied explanation; see Chapter04/SOURCES.md.
\subsubsection{প্রশ্নের চাহিদা থেকে ট্যাগ নির্বাচন}
চিত্র বা উদ্দীপক দেখে আগে তথ্যের সম্পর্ক শনাক্ত করবে। ধারাবাহিক ধাপে কাজ থাকলে ক্রমযুক্ত তালিকা, নাম-নম্বরের তুলনা থাকলে সারণি, অন্য পেজে যাওয়ার কাজ থাকলে লিংক এবং ব্যবহারকারীর তথ্য নেওয়ার কাজ থাকলে ফর্ম নির্বাচন করবে।

তারপর চাহিদা অনুযায়ী অ্যাট্রিবিউট বসাবে: লিংকের গন্তব্যে \texttt{href}, ছবির উৎসে \texttt{src} এবং বিস্তৃত ঘরে \texttt{rowspan}/\texttt{colspan}। উত্তর পরীক্ষা করার সময় প্রতিটি নির্বাচনের কারণ এক বাক্যে বলতে পারলে কোড মুখস্থ না করেও নতুন প্রশ্ন সমাধান করা যায়।
